% =========================================================
% Introduction and Geometry
% =========================================================

\section{Introduction}

\subsection{Motivation}

One of the central open problems in gravitational physics concerns the structural behavior of spacetime geometry in high-curvature vacuum regimes. In particular, classical Schwarzschild geometry predicts the formation of curvature singularities associated with geodesic incompleteness and divergent curvature invariants \cite{wald1984,hawkingellis1973,penrose1965,hawking1970}.

A variety of approaches have attempted to regulate such singular structures through modified geometric sectors, effective matter models, or quantum-gravitational corrections \cite{bardeen1968,dymnikova1992,hayward2006,ansoldi2008}. However, many such constructions either introduce uncontrolled ultraviolet assumptions or lack mathematically stable perturbative formulations.

The present work develops a conservative operator-theoretic perturbative framework associated with a representative bounded-curvature vacuum geometry compatible with the branch structure established in Topological Integrity Gravity II (TIG2).

The goal of the present paper is not to propose a complete alternative theory of gravity, but rather to define a mathematically coherent perturbative admissibility framework suitable for further spectral and geometric analysis.

% =========================================================

\subsection{TIG2 Structural Background}

The present framework builds upon the cubic branch structure introduced in TIG2 \cite{tig2}.

The central structural relation is

\begin{equation}
x^3 - x^2 + \beta^3 = 0,
\end{equation}

where

\begin{equation}
x = \frac{r_H}{2M},
\qquad
\beta = \frac{r_c}{2M}.
\end{equation}

Here:
\begin{itemize}[leftmargin=0.6cm]
\item $r_H$ denotes the horizon radius,
\item $M$ is the mass parameter,
\item and $r_c$ is a regulated representative core scale.
\end{itemize}

The discriminant structure is given by

\begin{equation}
\Delta(\beta)
=
\beta^3(4 - 27\beta^3).
\end{equation}

The critical branch-merger point occurs at

\begin{equation}
x_c = \frac{2}{3},
\qquad
\beta_c^3 = \frac{4}{27}.
\end{equation}

This branch structure defines the admissible near-critical sector underlying the TIG3 vacuum program.

% =========================================================

\subsection{Scientific Scope}

The present paper develops a conditional perturbative operator framework associated with a representative bounded-curvature geometry.

The analysis focuses on:
\begin{itemize}[leftmargin=0.6cm]
\item representative vacuum geometry,
\item perturbative admissibility,
\item operator-theoretic structure,
\item spectral consistency,
\item and bounded linear evolution.
\end{itemize}

The framework does not currently establish:
\begin{itemize}[leftmargin=0.6cm]
\item nonlinear stability,
\item complete covariant field equations,
\item quantum-gravity completion,
\item observational confirmation,
\item or full singularity resolution.
\end{itemize}

All perturbative statements are conditional upon the operator-theoretic assumptions explicitly stated below.

% =========================================================

\section{Representative TIG3 Geometry}

\subsection{Metric Ansatz}

We consider a static spherically symmetric representative geometry of the form

\begin{equation}
ds^2
=
-F(r)\,dt^2
+
\frac{dr^2}{F(r)}
+
r^2 d\Omega^2,
\end{equation}

where

\begin{equation}
d\Omega^2
=
d\theta^2
+
\sin^2\theta\, d\phi^2.
\end{equation}

The representative lapse function is chosen as

\begin{equation}
F(r)
=
1
-
\frac{2Mr^2}{r^3 + r_c^3}.
\label{eq:lapse}
\end{equation}

The geometry preserves Schwarzschild-compatible asymptotics while introducing a regulated high-curvature core scale.

% =========================================================

\subsection{Asymptotic Structure}

For large radial distance,

\begin{equation}
r \gg r_c,
\end{equation}

the lapse function admits the asymptotic expansion

\begin{equation}
F(r)
=
1
-
\frac{2M}{r}
+
\mathcal{O}\!\left(\frac{r_c^3}{r^3}\right).
\end{equation}

Hence the representative geometry reproduces Schwarzschild-compatible asymptotic behavior in the low-curvature regime \cite{wald1984,chandrasekhar1983}.

This compatibility condition is essential for maintaining consistency with the classical far-field sector of General Relativity.

% =========================================================

\subsection{Near-Core Structure}

Near the regulated core,

\begin{equation}
r \ll r_c,
\end{equation}

the lapse function behaves as

\begin{equation}
F(r)
\approx
1
-
\frac{2M}{r_c^3}r^2.
\end{equation}

The resulting geometry possesses bounded representative curvature invariants within the analyzed sector.

In particular, the representative Kretschmann scalar remains finite near the regulated core.

The present work interprets this structure conservatively as a bounded-curvature representative geometry rather than a complete singularity-resolution theorem.

% =========================================================

\subsection{Horizon Structure}

The horizon sector is determined by the roots of

\begin{equation}
F(r_H)=0.
\end{equation}

After rescaling according to

\begin{equation}
x
=
\frac{r_H}{2M},
\qquad
\beta
=
\frac{r_c}{2M},
\end{equation}

the horizon equation reduces exactly to the TIG2 cubic branch relation

\begin{equation}
x^3 - x^2 + \beta^3 = 0.
\end{equation}

Thus the representative TIG3 geometry is structurally compatible with the admissible branch hierarchy established in TIG2 \cite{tig2}.

% =========================================================

\subsection{Near-Critical Sector}

Near the critical branch-merger regime,

\begin{equation}
\beta \approx \beta_c,
\end{equation}

the geometry may develop elongated throat structures and modified perturbative trapping behavior.

The present work does not establish a rigorous near-critical resonance theorem. Such structures are interpreted only as conjectural perturbative sectors requiring further operator-theoretic analysis.

% =========================================================

\subsection{Interpretation}

The representative geometry introduced above should be interpreted as a mathematically controlled bounded-curvature sector suitable for perturbative operator analysis.

The present framework is intentionally conservative.

In particular, the geometry is not currently claimed to represent:
\begin{itemize}[leftmargin=0.6cm]
\item a complete covariant gravitational theory,
\item a final black-hole replacement model,
\item or a quantum-gravitational completion.
\end{itemize}

Rather, the geometry serves as the background structure for the perturbative admissibility program developed in the following sections.
